\documentclass[aspectratio=169]{beamer}
\usepackage{amsmath,amssymb}
\usepackage[T1]{fontenc}
\usepackage[utf8]{inputenc}
\usepackage{colortbl,xcolor,array}
\usepackage{tcolorbox}
\tcbuselibrary{skins}

% ── Colours ───────────────────────────────────────────────────────────
\definecolor{orange}{RGB}{230,126,34}
\definecolor{darkblue}{RGB}{44,62,80}
\definecolor{rowred}{RGB}{255,220,220}
\definecolor{rowblue}{RGB}{220,235,255}
% Wolfram palette (teal)
\definecolor{wfg}{RGB}{0,100,120}
\definecolor{wbg}{RGB}{228,246,250}
\definecolor{wfr}{RGB}{0,128,150}
% SymPy palette (green)
\definecolor{sfg}{RGB}{30,130,76}
\definecolor{sbg}{RGB}{232,246,237}
\definecolor{sfr}{RGB}{39,174,96}
% Setup palette (indigo)
\definecolor{setupcol}{RGB}{90,90,190}

% ── Beamer theme ──────────────────────────────────────────────────────
\usetheme{default}
\setbeamertemplate{navigation symbols}{}
\setbeamertemplate{itemize item}{\textcolor{orange}{\textbullet}}
\setbeamertemplate{itemize subitem}{\textcolor{orange}{\textbullet}}
\setbeamercolor{frametitle}{fg=orange,bg=white}
\setbeamercolor{normal text}{fg=darkblue}
\setbeamerfont{frametitle}{size=\LARGE,series=\bfseries}
\setbeamertemplate{footline}{%
  \hfill\footnotesize\color{gray}%
  \insertframenumber\,/\,\inserttotalframenumber\hspace{4pt}\vspace{4pt}%
}
\addtobeamertemplate{frametitle}{}%
  {\vspace{-4pt}\textcolor{orange!40}{\hrule height 0.4pt}}

% ── tcolorbox styles ──────────────────────────────────────────────────
\tcbset{
  wfstyle/.style={
    colback=wbg, colframe=wfr, colbacktitle=wfr, coltitle=white,
    fonttitle=\bfseries\tiny, title=Wolfram,
    fontupper=\ttfamily\scriptsize\color{wfg},
    boxrule=0.5pt, arc=3pt,
    top=3pt, bottom=3pt, left=5pt, right=5pt,
    toptitle=1pt, bottomtitle=1pt,
  },
  spstyle/.style={
    colback=sbg, colframe=sfr, colbacktitle=sfr, coltitle=white,
    fonttitle=\bfseries\tiny, title=SymPy,
    fontupper=\ttfamily\scriptsize\color{sfg},
    boxrule=0.5pt, arc=3pt,
    top=3pt, bottom=3pt, left=5pt, right=5pt,
    toptitle=1pt, bottomtitle=1pt,
  },
  setupstyle/.style={
    colback=setupcol!8, colframe=setupcol!50,
    colbacktitle=setupcol!70, coltitle=white,
    fonttitle=\bfseries\tiny, title=SymPy --- imports (run once),
    fontupper=\ttfamily\scriptsize\color{setupcol},
    boxrule=0.5pt, arc=3pt,
    top=2pt, bottom=2pt, left=5pt, right=5pt,
    toptitle=1pt, bottomtitle=1pt,
  }
}

% ── Inline helpers ────────────────────────────────────────────────────
\newcommand{\wcode}[1]{\texttt{\textcolor{wfg}{#1}}}
\newcommand{\scode}[1]{\texttt{\textcolor{sfg}{#1}}}
% Dual-code columns (usage: \dualcols{wolfram content}{sympy content})
\newcommand{\dualcols}[2]{%
  \begin{columns}[T,onlytextwidth]
    \begin{column}{0.487\textwidth}
      \begin{tcolorbox}[wfstyle]#1\end{tcolorbox}
    \end{column}
    \begin{column}{0.487\textwidth}
      \begin{tcolorbox}[spstyle]#2\end{tcolorbox}
    \end{column}
  \end{columns}
}

\title{Introduction to Discrete Mathematics}
\subtitle{Proposition Logic}
\date{}

% ══════════════════════════════════════════════════════════════════════
\begin{document}

% ── Title ─────────────────────────────────────────────────────────────
\begin{frame}[t]
  \vspace{3em}\centering
  {\usebeamerfont{title}\usebeamercolor[fg]{title}
    Introduction to Discrete Mathematics\par}
  \vspace{0.4em}
  {\large\textcolor{darkblue}{Proposition Logic}\par}
\end{frame}

% ── Overview ──────────────────────────────────────────────────────────
\begin{frame}[t]{Overview}
  \begin{itemize}\setlength\itemsep{0.6em}
    \item Proposition logic is a branch of logic studying ways of modifying and joining propositions.
    \item Propositions are crucial in mathematics, as all proofs rely on propositional logic.
    \item It was first seriously studied by Aristotle, who was also the first to popularize most topics covered in this section.
    \item You will learn about propositions and logical operators in this lesson.
  \end{itemize}
\end{frame}

% ── Proposition ───────────────────────────────────────────────────────
\begin{frame}[t]{Proposition}
  A \textbf{proposition} is a mathematical statement that is either \emph{true} or \emph{false}:
  \begin{itemize}\setlength\itemsep{0.05em}
    \item This course is on discrete mathematics. \quad
          \textbullet\quad 3 is greater than 4.
  \end{itemize}
  Interrogative, exclamatory or imperative sentences are \textbf{not} propositions:
  \begin{itemize}\setlength\itemsep{0.05em}
    \item Why are some sheep black? \quad
          \textbullet\quad Do the dishes.
  \end{itemize}
  \vspace{0.3em}
  \begin{tcolorbox}[setupstyle]
    from sympy import symbols, Not, And, Or, Implies, Equivalent,
    simplify\\
    from sympy.logic.boolalg import truth\_table\\
    p, q = symbols('p q')
  \end{tcolorbox}
  \vspace{0.2em}
  \dualcols{%
    In[1]:= \{Boole[True], Boole[False]\}\\
    Out[1]= \{1, 0\}
  }{%
    [int(True), int(False)] \# [1, 0]
  }
\end{frame}

% ── Basic Operators: NOT + AND ────────────────────────────────────────
\begin{frame}[t]{Basic Operators}
  Let $p$ and $q$ be propositions.
  \smallskip

  The \textbf{negation} of $p$
  \,($\neg$, \wcode{!}, \wcode{Not})\,
  states ``$p$ is false'':
  \vspace{2pt}
  \dualcols{%
    \{$\neg$p,\ !p,\ Not[p]\}
  }{%
    Not(p)\quad\# or \textasciitilde{}p
  }

  \begin{center}
    \renewcommand{\arraystretch}{1.2}
    \begin{tabular}{|>{\centering}p{1.5cm}|>{\centering\arraybackslash}p{1.5cm}|}
      \hline $p$ & $\neg p$\\\hline
      \rowcolor{rowred}  True  & False\\
      \rowcolor{rowblue} False & True\\\hline
    \end{tabular}
  \end{center}
  \vspace{2pt}

  The \textbf{conjunction} of $p$ and $q$
  \,($\wedge$, \wcode{\&\&}, \wcode{And})\,
  is ``$p$ and $q$'':
  \vspace{2pt}
  \dualcols{%
    \{p $\wedge$ q,\ p \&\& q,\ And[p,q]\}
  }{%
    And(p, q)\quad\# or p \& q
  }
\end{frame}

% ── Basic Operators: OR ───────────────────────────────────────────────
\begin{frame}[t]{Basic Operators (continued)}
  The \textbf{disjunction} of $p$ and $q$
  \,($\vee$, \wcode{||}, \wcode{Or})\,
  is ``$p$ or $q$'':
  \vspace{4pt}
  \dualcols{%
    \{p $\vee$ q,\ p || q,\ Or[p,q]\}
  }{%
    Or(p, q)\quad\# or p | q
  }

  \vspace{0.5em}
  \begin{center}
    \renewcommand{\arraystretch}{1.25}
    \begin{tabular}{|>{\centering}p{1.5cm}|>{\centering}p{1.5cm}|>{\centering\arraybackslash}p{1.8cm}|}
      \hline $p$ & $q$ & $p \vee q$\\\hline
      \rowcolor{rowred}  True  & True  & True\\
      \rowcolor{rowred}  True  & False & True\\
      \rowcolor{rowred}  False & True  & True\\
      \rowcolor{rowblue} False & False & False\\\hline
    \end{tabular}
  \end{center}

  \vspace{0.5em}
  \begin{tcolorbox}[spstyle, title=SymPy --- truth table]
    for row in truth\_table(Or(p, q), [p, q]):\\
    \quad\quad print(row)\quad\# prints (vals, result) tuples
  \end{tcolorbox}
\end{frame}

% ── Example 1 ─────────────────────────────────────────────────────────
\begin{frame}[t]{Example}
  Let $p$ = ``I studied for the exam'' and $q$ = ``I passed the exam.''
  \vspace{0.5em}

  \textbf{Express in English:}
  \begin{itemize}
    \item $\neg q$ \hfill I did not pass the exam.
    \item $p \wedge q$ \hfill I studied for the exam and I passed the exam.
  \end{itemize}
  \vspace{0.5em}

  \textbf{Express in logic:}
  \begin{itemize}
    \item I did not pass the exam, nor did I study for it.
          \hfill $\neg q \wedge \neg p$
    \item I studied for the exam, but I did not pass it.
          \hfill $p \wedge \neg q$
  \end{itemize}
  \vspace{0.5em}

  \begin{tcolorbox}[spstyle, title=SymPy --- verify]
    Not(q)\quad\quad\quad\quad\quad\quad\quad\quad
    \# \textasciitilde{}q\\
    And(p, q)\quad\quad\quad\quad\quad\quad\quad
    \# p \& q\\
    And(Not(q), Not(p))\quad\quad\# \textasciitilde{}q \& \textasciitilde{}p\\
    And(p, Not(q))\quad\quad\quad\quad\# p \& \textasciitilde{}q
  \end{tcolorbox}
\end{frame}

% ── More Operators: IMPLIES + BICONDITIONAL ───────────────────────────
\begin{frame}[t]{More Operators}
  The \textbf{conditional} $p \Rightarrow q$
  \,($\Rightarrow$, \wcode{Implies})\,
  is ``If $p$, then $q$'':
  \vspace{2pt}
  \dualcols{%
    \{p $\Rightarrow$ q,\ Implies[p,q]\}
  }{%
    Implies(p, q)
  }

  \begin{center}
    \renewcommand{\arraystretch}{1.2}
    \begin{tabular}{|>{\centering}p{1.5cm}|>{\centering}p{1.5cm}|>{\centering\arraybackslash}p{2cm}|}
      \hline $p$ & $q$ & $p \Rightarrow q$\\\hline
      \rowcolor{rowred}  True  & True  & True\\
      \rowcolor{rowblue} True  & False & False\\
      \rowcolor{rowred}  False & True  & True\\
      \rowcolor{rowred}  False & False & True\\\hline
    \end{tabular}
  \end{center}
  \vspace{2pt}

  The \textbf{biconditional} $p \Leftrightarrow q$
  \,($\Leftrightarrow$, \wcode{Equivalent})\,
  is ``$p$ iff $q$'':
  \vspace{2pt}
  \dualcols{%
    \{p $\Leftrightarrow$ q,\ Equivalent[p,q]\}
  }{%
    Equivalent(p, q)
  }
\end{frame}

% ── Example 2 ─────────────────────────────────────────────────────────
\begin{frame}[t]{Example}
  Let $p$ = ``It is raining'' and $q$ = ``Kala wants to play outside.''
  \vspace{0.5em}

  \textbf{Express in English:}
  \begin{itemize}
    \item $p \Rightarrow q$ \hfill
          If it is raining, then Kala wants to play outside.
    \item $\neg q \Leftrightarrow p$ \hfill
          Kala does not want to play outside iff it is raining.
  \end{itemize}
  \vspace{0.5em}

  \textbf{Express in logic:}
  \begin{itemize}
    \item Either Kala wants to play outside or it is raining.
          \hfill $q \vee p$
    \item Kala wants to play outside regardless of rain.
          \hfill $(p \Rightarrow q) \wedge (\neg p \Rightarrow q)$
  \end{itemize}
  \vspace{0.4em}

  \begin{tcolorbox}[spstyle, title=SymPy --- simplification]
    expr = And(Implies(p,q), Implies(Not(p),q))\\
    simplify(expr)\quad\# simplifies to just: q
  \end{tcolorbox}
\end{frame}

% ── Precedence ────────────────────────────────────────────────────────
\begin{frame}[t]{Precedence}
  The precedence of logical operators is as follows:
  \begin{center}
    \large$() \;>\; \neg \;>\; \wedge \;>\; \vee \;>\; \Rightarrow \;>\; \Leftrightarrow$
  \end{center}
  \begin{center}
    \colorbox{orange!20}{\textbf{Mnemonic: P\,N\,A\,O\,I\,B}
    (Parentheses, Not, And, Or, Implies, Biconditional)}
  \end{center}
  \vspace{0.4em}
  \textbf{Evaluate if $p$ is true and $q$ is false:}
  \vspace{2pt}
  \dualcols{%
    In[2]:= \{p,q\} = \{True,False\}\\
    In[3]:= $\neg$p $\vee$ q $\Rightarrow$ p $\wedge$ $\neg$q\\
    Out[3]= True
  }{%
    p\_v, q\_v = True, False\\
    e = Implies(\textasciitilde{}p|q, p\&\textasciitilde{}q)\\
    e.subs(\{p:p\_v, q:q\_v\})\quad\# True
  }
  \vspace{0.3em}
  \dualcols{%
    $\neg$p $\vee$ q $\vee$ (p $\wedge$ q $\vee$ True) $\Rightarrow$ p $\wedge$ q
  }{%
    f = Implies(\textasciitilde{}p|q|(p\&q)|True, p\&q)\\
    f.subs(\{p:p\_v, q:q\_v\})\quad\# False
  }
\end{frame}

% ── Converse / Contrapositive ─────────────────────────────────────────
\begin{frame}[t]{Converse, Contrapositive, Inverse or Opposite?}
  Consider the conditional proposition $p \Rightarrow q$.
  \vspace{0.4em}

  \begin{columns}[T,onlytextwidth]
    \begin{column}{0.44\textwidth}
      \begin{itemize}\setlength\itemsep{0.45em}
        \item \textbf{Converse:}\quad $q \Rightarrow p$
        \item \textbf{Contrapositive:}\quad $\neg q \Rightarrow \neg p$
        \item \textbf{Inverse:}\quad $\neg p \Rightarrow \neg q$
        \item \textbf{Opposite:}\quad $\neg(p \Rightarrow q)$
      \end{itemize}
      \vspace{0.6em}
      The \textbf{contrapositive} will be useful in future lessons.
    \end{column}
    \begin{column}{0.53\textwidth}
      \begin{tcolorbox}[spstyle]
        converse = Implies(q, p)\\
        contra   = Implies(\textasciitilde{}q, \textasciitilde{}p)\\
        inverse  = Implies(\textasciitilde{}p, \textasciitilde{}q)\\
        opposite = Not(Implies(p, q))\\[4pt]
        \# Contrapositive $\equiv$ original:\\
        Equivalent(Implies(p,q),\\
        \quad\quad\quad\quad Implies(\textasciitilde{}q,\textasciitilde{}p))
        \# True
      \end{tcolorbox}
    \end{column}
  \end{columns}
\end{frame}

% ── Logic Gates ───────────────────────────────────────────────────────
\begin{frame}[t]{Logic Gates}
  A \textbf{logic gate} is a simple switching circuit with an output corresponding to logical operators.
  \vspace{0.5em}

  Logic gates are usually represented in diagrams:
  \vspace{0.4em}

  \begin{center}
    \renewcommand{\arraystretch}{1.6}
    \begin{tabular}{ccc}
      \textbf{NOT gate} & \textbf{AND gate} & \textbf{OR gate}\\[4pt]
      $p \longrightarrow \framebox{$\neg$} \longrightarrow \neg p$
      &
      $\begin{array}{r}p\\q\end{array}
       \longrightarrow \framebox{$\wedge$} \longrightarrow p\wedge q$
      &
      $\begin{array}{r}p\\q\end{array}
       \longrightarrow \framebox{$\vee$} \longrightarrow p\vee q$
    \end{tabular}
  \end{center}

  \vspace{0.5em}
  \begin{tcolorbox}[spstyle, title=SymPy --- gate truth tables]
    for row in truth\_table(Not(p), [p]):\\
    \quad\quad print(row)\quad\# NOT\\
    for row in truth\_table(And(p,q), [p,q]):\\
    \quad\quad print(row)\quad\# AND\\
    for row in truth\_table(Or(p,q), [p,q]):\\
    \quad\quad print(row)\quad\# OR
  \end{tcolorbox}

  They are essential in understanding computing \textbf{from the ground up}.
\end{frame}

% ── Summary ───────────────────────────────────────────────────────────
\begin{frame}[t]{Summary}
  \begin{itemize}\setlength\itemsep{0.6em}
    \item A \textbf{proposition} is a mathematical statement that is either true or false, but never both.
    \item The three main logical operators are
          \textbf{negation} ($\neg$), \textbf{conjunction} ($\wedge$), and \textbf{disjunction} ($\vee$).
    \item Many other operators exist, most importantly the
          \textbf{conditional} ($\Rightarrow$) and \textbf{biconditional} ($\Leftrightarrow$).
    \item Logical expressions easily get convoluted.
  \end{itemize}
  \vspace{0.8em}
  \begin{tcolorbox}[spstyle, title=SymPy --- key imports for this lesson]
    from sympy import symbols, Not, And, Or, Implies, Equivalent,
    simplify\\
    from sympy.logic.boolalg import truth\_table, to\_cnf, to\_dnf\\
    from sympy.logic.inference import satisfiable\\
    p, q = symbols('p q')
  \end{tcolorbox}
\end{frame}

\end{document}
